\documentclass[12pt]{article}

\usepackage{amsmath, amssymb}
\usepackage{geometry}
\usepackage{titlesec}
\usepackage{mdframed}
\usepackage{xcolor}

\geometry{a4paper, top=2cm, bottom=2cm, left=2.5cm, right=2.5cm}

\setlength{\parskip}{4pt}
\setlength{\parindent}{0pt}
\setlength{\abovedisplayskip}{6pt}
\setlength{\belowdisplayskip}{6pt}
\setlength{\abovedisplayshortskip}{3pt}
\setlength{\belowdisplayshortskip}{3pt}

\definecolor{noteblue}{RGB}{25, 80, 150}
\definecolor{boxgray}{RGB}{246, 248, 252}

\titleformat{\section}{\large\bfseries}{}{0em}{}[\vspace{2pt}\titlerule\vspace{4pt}]
\titlespacing{\section}{0pt}{10pt}{4pt}
\titleformat{\subsection}{\normalsize\bfseries}{}{0em}{}
\titlespacing{\subsection}{0pt}{8pt}{2pt}

\begin{document}

\begin{center}
  {\LARGE\bfseries Daily JEE Advanced Problems}\\[4pt]
  {\large JEE Advanced 2025 --- Paper I}\\[2pt]
  {\small\color{gray} Daily Problem Series \quad|\quad Virat Dixit}
\end{center}

\vspace{4pt}
\noindent\rule{\linewidth}{0.8pt}
\vspace{4pt}

\section*{Problem Statement}

Let $\mathbb{R}$ denote the set of all real numbers. Let $f:\mathbb{R}\to\mathbb{R}$ be a function
such that $f(x)>0$ for all $x\in\mathbb{R}$, and
\[
  f(x+y)=f(x)\,f(y)\quad\text{for all }x,y\in\mathbb{R}.
\]
Let the real numbers $a_1,a_2,a_3,\ldots$ with $a_1>0$ be an A.P.
If $f(a_{31})=64\cdot f(a_{25})$ and
\[
  \sum_{i=1}^{50}f(a_i)=3\!\left(2^{25}+1\right),
\]
find the value of $\displaystyle\sum_{i=6}^{30}f(a_i)$.

\vspace{6pt}
\begin{mdframed}[backgroundcolor=boxgray, linecolor=noteblue, linewidth=0.8pt,
                 innertopmargin=5pt, innerbottommargin=5pt,
                 innerleftmargin=8pt, innerrightmargin=8pt,
                 skipabove=4pt, skipbelow=4pt]
\textbf{\color{noteblue}Note.} This problem combines functional equations, arithmetic progressions,
geometric progressions, factorisation identities, and symmetry in summation.
\end{mdframed}

\section*{Solution}

\subsection*{Step 1: Identify the form of $f$}

Since $f(x+y)=f(x)f(y)$ with $f(x)>0$, the function must be exponential:
$f(x)=e^{kx}$ for some real constant $k$.

\subsection*{Step 2: Express $f(a_n)$ as a G.P.}

Let the A.P.\ have first term $a_1$ and common difference $d$, so $a_n=a_1+(n-1)d$. Then
\[
  f(a_n)=f\!\bigl(a_1+(n-1)d\bigr)=f(a_1)\cdot[f(d)]^{n-1}.
\]
Setting $A=f(a_1)>0$ and $r=f(d)>0$:
\[
  f(a_n)=A\cdot r^{n-1}. \tag{1}
\]
So $f(a_1),f(a_2),\ldots,f(a_{50})$ form a \textbf{G.P.} with first term $A$ and common ratio $r$.

\subsection*{Step 3: Find $r$}

From $f(a_{31})=64\,f(a_{25})$ and (1):
\[
  Ar^{30}=64\cdot Ar^{24}
  \implies r^6=64=2^6
  \implies \boxed{r=2.}
\]

\subsection*{Step 4: Find $A$}

Sum of first 50 terms:
\[
  \sum_{i=1}^{50}f(a_i)=A(1+2+\cdots+2^{49})=A(2^{50}-1).
\]
Given this equals $3(2^{25}+1)$, and factorising $2^{50}-1=(2^{25}+1)(2^{25}-1)$:
\[
  A(2^{25}+1)(2^{25}-1)=3(2^{25}+1)
  \implies A(2^{25}-1)=3
  \implies \boxed{A=\dfrac{3}{2^{25}-1}.} \tag{2}
\]

\subsection*{Step 5: Compute the required sum}

\[
  \sum_{i=6}^{30}f(a_i)
  =A(2^5+2^6+\cdots+2^{29})
  =A\cdot 2^5(2^{25}-1).
\]
Substituting $A$ from (2):
\[
  \sum_{i=6}^{30}f(a_i)
  =\frac{3}{2^{25}-1}\cdot 32\cdot(2^{25}-1)
  =3\times 32
  =\boxed{96.}
\]

\vspace{8pt}
\noindent\rule{\linewidth}{0.4pt}
\vspace{4pt}
\begin{center}
\itshape
``Problems worthy of attack prove their worth by fighting back.''\\[1pt]
--- Paul Erd\H{o}s
\end{center}

\end{document}
